\documentclass[12pt]{article}
\usepackage{amsmath,amssymb,amsthm}
\usepackage{geometry}
\geometry{margin=1in}
\title{Unifying Stability in PDEs: A Single Invariance Perspective}
\author{}
\date{}

% Theorem and Lemma environments
\newtheorem{theorem}{Theorem}
\newtheorem{lemma}{Lemma}
\newtheorem{definition}{Definition}

\begin{document}

\maketitle

\begin{abstract}
This document unifies the principles of energy boundedness, spectral decay, compactness, and geometric gradient control into a single invariance framework for partial differential equations (PDEs). By interpreting PDE dynamics as gradient flows in suitable Hilbert spaces, we show that the so-called "four pillars" of stability are not independent processes but facets of a single residual functional \(\mathcal{R}(u)\). The framework is extended to bounded domains and compressible flows, and links between spectral suppression and geometric curvature are explored rigorously. This approach provides a universal lens for understanding PDE stability across various contexts.
\end{abstract}

\section{Introduction}

The stability and regularity of PDEs, such as the 3D Navier--Stokes equations, have traditionally been analyzed through four key mechanisms:
\begin{enumerate}
    \item \textbf{Energy Boundedness:} Ensuring total energy and enstrophy remain finite.
    \item \textbf{Spectral Decay:} Suppressing high-frequency (small-scale) energy cascades.
    \item \textbf{Compactness:} Establishing strong convergence in appropriate function spaces.
    \item \textbf{Gradient Control:} Damping local gradients to prevent singularities.
\end{enumerate}
These mechanisms are often treated as distinct steps in a feedback loop. However, this decomposition is a methodological artifact. In reality, the PDE enforces a single invariance condition that simultaneously satisfies all these aspects. This document formalizes this perspective and extends it to bounded domains, compressible flows, and spectral-geometric connections.

\section{Formalizing the Single Invariance Perspective}

\subsection{Master Functional \(\Phi(u)\)}
We propose that all stability constraints are encoded in a single functional \(\Phi(u)\), which captures:
\begin{itemize}
    \item Total energy: \(E(u) = \frac{1}{2} \| u \|_{L^2}^2\).
    \item Dissipation: \(D(u) = \frac{\nu}{2} \| \nabla u \|_{L^2}^2\).
    \item Spectral suppression: \(S(u) = \sup_{k > k_0} \frac{E(k)}{k^\beta}\), where \(E(k)\) is the spectral energy density.
    \item Forcing: \(F(u) = -\int f \cdot u \, dx\).
\end{itemize}
Thus, the functional becomes:
\begin{equation}
\Phi(u) = E(u) + D(u) + S(u) + F(u).
\end{equation}

\subsection{Residual Functional \(\mathcal{R}(u)\)}
The PDE is expressed as the stationarity condition of \(\Phi(u)\):
\begin{equation}
\mathcal{R}(u) = \nabla \Phi(u).
\end{equation}
For the Navier--Stokes equations, this becomes:
\begin{equation}
\mathcal{R}(u) = -\nu \Delta u + (u \cdot \nabla) u + \nabla p - f.
\end{equation}

\subsection{Proving That \(\mathcal{R}(u) = 0\) Implies All Constraints}

\paragraph{Energy Boundedness:} From \(\Phi(u)\), monotonicity ensures bounded energy:
\begin{equation}
\frac{d}{dt} \Phi(u) = -\|\nabla \Phi(u)\|^2 \leq 0 \implies \|u\|_{L^2}^2 \text{ remains finite.}
\end{equation}

\paragraph{Spectral Decay:} The term \(\sup_{k > k_0} \frac{E(k)}{k^\beta}\) suppresses high-frequency modes, ensuring:
\begin{equation}
E(k) \leq C k^{-\beta}, \quad \beta > 3.
\end{equation}

\paragraph{Compactness:} Strong convergence in \(L^2\) is ensured by the Aubin--Lions lemma and spectral decay:
\begin{equation}
\sup_{k > k_0} \frac{E(k)}{k^\beta} \to 0 \quad \text{as } k_0 \to \infty.
\end{equation}

\paragraph{Gradient Control:} Dissipation directly bounds local gradients:
\begin{equation}
\|\nabla u\|_{L^2}^2 \leq \frac{\Phi(u)}{\nu}.
\end{equation}

\section{Extending \(\Phi(u)\) to Bounded Domains or Compressible Flows}

\subsection{Bounded Domains}
For \(\Omega \subset \mathbb{R}^3\), include boundary terms:
\begin{equation}
\Phi_B(u) = \int_{\partial \Omega} \mathcal{B}(u) \, dS, \quad \mathcal{B}(u) = \nu (\nabla u \cdot n) u.
\end{equation}
The modified functional is:
\begin{equation}
\Phi(u) = \Phi_E(u) + \Phi_D(u) + S(u) + \Phi_B(u) - \int_\Omega f \cdot u \, dx.
\end{equation}

\subsection{Compressible Flows}
For compressible Navier--Stokes:
\begin{equation}
\Phi(\rho, u) = \frac{1}{2} \int_\Omega \rho |u|^2 \, dx + \int_\Omega P(\rho) \, dx + \Phi_D(u) + S(u) - \int_\Omega f \cdot u \, dx.
\end{equation}
Here, \(P(\rho)\) is the internal pressure potential.

\section{Spectral-Geometric Connections for Curvature-Based Damping}

Positive Ricci curvature \(\text{Ric}(g) \geq \kappa > 0\) ensures spectral gaps for the Laplace--Beltrami operator:
\begin{equation}
\lambda_k \geq C_1 k^2 + C_2 \text{Ric}.
\end{equation}
Spectral suppression follows:
\begin{equation}
E(k) \sim \lambda_k^{-\beta}, \quad \beta > 3.
\end{equation}

\section{Conclusion}
We have unified the principles of PDE stability into a single invariance condition \(\mathcal{R}(u) = 0\), with all constraints encoded in \(\Phi(u)\). Extensions to bounded domains, compressible flows, and spectral-geometric connections demonstrate the broad applicability of this approach. Future work includes further refinements for boundary effects and rigorous compactness proofs in non-ideal geometries.

\end{document}
